\documentclass[10pt,a4paper]{amsart}
\usepackage[margin=19mm]{geometry}
\usepackage{amsmath,amssymb,mathtools}
\usepackage[scr=rsfs,cal=euler]{mathalfa}
\usepackage{xfrac}
\usepackage[unicode,hidelinks,bookmarks=false]{hyperref}
\newcommand{\ie}{i.e.\ }
\newcommand{\cf}{cf.\ }
\newcommand{\resp}{respectively\ }
\newcommand{\Subsection}{\subsection}
\providecommand{\nobreakdash}{\nobreak\hbox{-}}
\setlength{\emergencystretch}{2em}
\multlinegap0pt
% ---------------------------------------------------------------------------
% Editorial AMS wrapper prelude. The theorem environments and the mathematical macros below
% are transcribed verbatim from the paper's own preamble (cached-originals/source0.tex lines
% 62--359); the AMS amsart class, the named format packages of the pinned TeX Live runtime and
% this editorial formatting are not part of the author's text. No mathematics is changed.
% ---------------------------------------------------------------------------
\usepackage{enumitem}
\newtheorem{maintheorem}{Theorem}
\newtheorem{maincorollary}[maintheorem]{Corollary}
\renewcommand*{\themaintheorem}{\Alph{maintheorem}}
\renewcommand*{\themaincorollary}{\Alph{maincorollary}}
\newtheorem{theorem}{Theorem}[section]
\newtheorem{corollary}[theorem]{Corollary}
\newtheorem{lemma}[theorem]{Lemma}
\newtheorem{proposition}[theorem]{Proposition}
\newtheorem{problem}[theorem]{Problem}
\newtheorem{innercustomthm}{Theorem}
\newenvironment{customthm}[1]
{\renewcommand\theinnercustomthm{#1}\innercustomthm}
{\endinnercustomthm}
\theoremstyle{definition}
\newtheorem{definition}[theorem]{Definition}
\newtheorem{example}[theorem]{Example}
\theoremstyle{definition}
\newtheorem{remark}[theorem]{Remark}


\numberwithin{equation}{section}
\newcommand{\nrm}[1]{\Vert#1\Vert}
\newcommand{\abs}[1]{\vert#1\vert}
\newcommand{\brk}[1]{\langle#1\rangle}
\newcommand{\set}[1]{\{#1\}}
\newcommand{\tld}[1]{\widetilde{#1}}
\newcommand{\br}[1]{\overline{#1}}
\newcommand{\ubr}[1]{\underline{#1}}
\newcommand{\wht}[1]{\widehat{#1}}
\newcommand{\nnrm}[1]{{\vert\kern-0.25ex\vert\kern-0.25ex\vert #1 
		\vert\kern-0.25ex\vert\kern-0.25ex\vert}}

\newcommand{\dist}{\mathrm{dist}\,}
\newcommand{\sgn}{{\mathrm{sgn}}\,}
\newcommand{\supp}{{\mathrm{supp}}\,}
\newcommand{\tr}{\mathrm{tr}\,}

\renewcommand{\Re}{\mathrm{Re}}
\renewcommand{\Im}{\mathrm{Im}}

\newcommand{\aeq}{\simeq}
\newcommand{\aleq}{\lesssim}
\newcommand{\ageq}{\gtrsim}

\newcommand{\sbeq}{\subseteq}
\newcommand{\speq}{\supseteq}

\newcommand{\lap}{\Delta}
\newcommand{\bx}{\Box}

\newcommand{\ud}{\mathrm{d}}
\newcommand{\rd}{\partial}
\newcommand{\nb}{\nabla}

\newcommand{\imp}{\Rightarrow}
\newcommand{\pmi}{\Leftarrow}
\newcommand{\impmi}{\Leftrightarrow}

\newcommand{\bb}{\Big}

\newcommand{\0}{\emptyset}
\newcommand{\hp}{{\bbR^2_+}}
\newcommand{\ohp}{\overline{\bbR^2_+}}

\newcommand{\alp}{\alpha}
\newcommand{\bt}{\beta}
\newcommand{\gmm}{\gamma}
\newcommand{\Gmm}{\Gamma}
\newcommand{\dlt}{\delta}
\newcommand{\Dlt}{\Delta}
\newcommand{\eps}{\epsilon}
\newcommand{\veps}{\varepsilon}
\newcommand{\kpp}{\kappa}
\newcommand{\lmb}{\kappa}
\newcommand{\Lmb}{\kappa}
\newcommand{\sgm}{\sigma}
\newcommand{\Sgm}{\Sigma}
\newcommand{\tht}{\theta}
\newcommand{\Tht}{\Theta}
\newcommand{\vtht}{\vartheta}
\newcommand{\omg}{\omega}
\newcommand{\Omg}{\Omega}
\newcommand{\ups}{\upsilon}
\newcommand{\zt}{\zeta}
\newcommand{\Zt}{\Zeta}

\newcommand{\bfa}{{\bf a}}
\newcommand{\bfb}{{\bf b}}
\newcommand{\bfc}{{\bf c}}
\newcommand{\bfd}{{\bf d}}
\newcommand{\bfe}{{\bf e}}
\newcommand{\bff}{{\bf f}}
\newcommand{\bfg}{{\bf g}}
\newcommand{\bfh}{{\bf h}}
\newcommand{\bfi}{{\bf i}}
\newcommand{\bfj}{{\bf j}}
\newcommand{\bfk}{{\bf k}}
\newcommand{\bfl}{{\bf l}}
\newcommand{\bfm}{{\bf m}}
\newcommand{\bfn}{{\bf n}}
\newcommand{\bfo}{{\bf o}}
\newcommand{\bfp}{{\bf p}}
\newcommand{\bfq}{{\bf q}}
\newcommand{\bfr}{{\bf r}}
\newcommand{\bfs}{{\bf s}}
\newcommand{\bft}{{\bf t}}
\newcommand{\bfu}{{\bf u}}
\newcommand{\bfv}{{\bf v}}
\newcommand{\bfw}{{\bf w}}
\newcommand{\bfx}{{\bf x}}
\newcommand{\bfy}{{\bf y}}
\newcommand{\bfz}{{\bf z}}
\newcommand{\bfrho}{{\boldsymbol{\rho}}}

\newcommand{\bfA}{{\bf A}}
\newcommand{\bfB}{{\bf B}}
\newcommand{\bfC}{{\bf C}}
\newcommand{\bfD}{{\bf D}}
\newcommand{\bfE}{{\bf E}}
\newcommand{\bfF}{{\bf F}}
\newcommand{\bfG}{{\bf G}}
\newcommand{\bfH}{{\bf H}}
\newcommand{\bfI}{{\bf I}}
\newcommand{\bfJ}{{\bf J}}
\newcommand{\bfK}{{\bf K}}
\newcommand{\bfL}{{\bf L}}
\newcommand{\bfM}{{\bf M}}
\newcommand{\bfN}{{\bf N}}
\newcommand{\bfO}{{\bf O}}
\newcommand{\bfP}{{\bf P}}
\newcommand{\bfQ}{{\bf Q}}
\newcommand{\bfR}{{\bf R}}
\newcommand{\bfS}{{\bf S}}
\newcommand{\bfT}{{\bf T}}
\newcommand{\bfU}{{\bf U}}
\newcommand{\bfV}{{\bf V}}
\newcommand{\bfW}{{\bf W}}
\newcommand{\bfX}{{\bf X}}
\newcommand{\bfY}{{\bf Y}}
\newcommand{\bfZ}{{\bf Z}}

\newcommand{\bfPi}{{\mathbf{\Pi}}}

\newcommand{\bbA}{\mathbb A}
\newcommand{\bbB}{\mathbb B}
\newcommand{\bbC}{\mathbb C}
\newcommand{\bbD}{\mathbb D}
\newcommand{\bbE}{\mathbb E}
\newcommand{\bbF}{\mathbb F}
\newcommand{\bbG}{\mathbb G}
\newcommand{\bbH}{\mathbb H}
\newcommand{\bbI}{\mathbb I}
\newcommand{\bbJ}{\mathbb J}
\newcommand{\bbK}{\mathbb K}
\newcommand{\bbL}{\mathbb L}
\newcommand{\bbM}{\mathbb M}
\newcommand{\bbN}{\mathbb N}
\newcommand{\bbO}{\mathbb O}
\newcommand{\bbP}{\mathbb P}
\newcommand{\bbQ}{\mathbb Q}
\newcommand{\bbR}{\mathbb R}
\newcommand{\bbS}{\mathbb S}
\newcommand{\bbT}{\mathbb T}
\newcommand{\bbU}{\mathbb U}
\newcommand{\bbV}{\mathbb V}
\newcommand{\bbW}{\mathbb W}
\newcommand{\bbX}{\mathbb X}
\newcommand{\bbY}{\mathbb Y}
\newcommand{\bbZ}{\mathbb Z}

\newcommand{\calA}{\mathcal A}
\newcommand{\calB}{\mathcal B}
\newcommand{\calC}{\mathcal C}
\newcommand{\calD}{\mathcal D}
\newcommand{\calE}{\mathcal E}
\newcommand{\calF}{\mathcal F}
\newcommand{\calG}{\mathcal G}
\newcommand{\calH}{\mathcal H}
\newcommand{\calI}{\mathcal I}
\newcommand{\calJ}{\mathcal J}
\newcommand{\calK}{\mathcal K}
\newcommand{\calL}{\mathcal L}
\newcommand{\calM}{\mathcal M}
\newcommand{\calN}{\mathcal N}
\newcommand{\calO}{\mathcal O}
\newcommand{\calP}{\mathcal P}
\newcommand{\calQ}{\mathcal Q}
\newcommand{\calR}{\mathcal R}
\newcommand{\calS}{\mathcal S}
\newcommand{\calT}{\mathcal T}
\newcommand{\calU}{\mathcal U}
\newcommand{\calV}{\mathcal V}
\newcommand{\calW}{\mathcal W}
\newcommand{\calX}{\mathcal X}
\newcommand{\calY}{\mathcal Y}
\newcommand{\calZ}{\mathcal Z}

\newcommand{\varep}{\varepsilon}
\newcommand{\dd}{\mathrm{d}}

\newcommand{\frkA}{\mathfrak A}
\newcommand{\frkB}{\mathfrak B}
\newcommand{\frkC}{\mathfrak C}
\newcommand{\frkD}{\mathfrak D}
\newcommand{\frkE}{\mathfrak E}
\newcommand{\frkF}{\mathfrak F}
\newcommand{\frkG}{\mathfrak G}
\newcommand{\frkH}{\mathfrak H}
\newcommand{\frkI}{\mathfrak I}
\newcommand{\frkJ}{\mathfrak J}
\newcommand{\frkK}{\mathfrak K}
\newcommand{\frkL}{\mathfrak L}
\newcommand{\frkM}{\mathfrak M}
\newcommand{\frkN}{\mathfrak N}
\newcommand{\frkO}{\mathfrak O}
\newcommand{\frkP}{\mathfrak P}
\newcommand{\frkQ}{\mathfrak Q}
\newcommand{\frkR}{\mathfrak R}
\newcommand{\frkS}{\mathfrak S}
\newcommand{\frkT}{\mathfrak T}
\newcommand{\frkU}{\mathfrak U}
\newcommand{\frkV}{\mathfrak V}
\newcommand{\frkW}{\mathfrak W}
\newcommand{\frkX}{\mathfrak X}
\newcommand{\frkY}{\mathfrak Y}
\newcommand{\frkZ}{\mathfrak Z}

\newcommand{\frka}{\mathfrak a}
\newcommand{\frkb}{\mathfrak b}
\newcommand{\frkc}{\mathfrak c}
\newcommand{\frkd}{\mathfrak d}
\newcommand{\frke}{\mathfrak e}
\newcommand{\frkf}{\mathfrak f}
\newcommand{\frkg}{\mathfrak g}
\newcommand{\frkh}{\mathfrak h}
\newcommand{\frki}{\mathfrak i}
\newcommand{\frkj}{\mathfrak j}
\newcommand{\frkk}{\mathfrak k}
\newcommand{\frkl}{\mathfrak l}
\newcommand{\frkm}{\mathfrak m}
\newcommand{\frkn}{\mathfrak n}
\newcommand{\frko}{\mathfrak o}
\newcommand{\frkp}{\mathfrak p}
\newcommand{\frkq}{\mathfrak q}
\newcommand{\frkr}{\mathfrak r}
\newcommand{\frks}{\mathfrak s}
\newcommand{\frkt}{\mathfrak t}
\newcommand{\frku}{\mathfrak u}
\newcommand{\frkv}{\mathfrak v}
\newcommand{\frkw}{\mathfrak w}
\newcommand{\frkx}{\mathfrak x}
\newcommand{\frky}{\mathfrak y}
\newcommand{\frkz}{\mathfrak z}
\newcommand{\To}{\longrightarrow}
\newcommand{\weakto}{\rightharpoonup}
\newcommand{\embed}{\righthookarrow}
\newcommand{\vect}[3]{\begin{pmatrix} #1 \\ #2 \\ #3 \end{pmatrix}}
\newcommand{\pfstep}[1]{\vskip.5em \noindent {\it #1.}}
\newcommand{\ackn}[1]{
	\addtocontents{toc}{\protect\setcounter{tocdepth}{1}}
	\subsection*{Acknowledgements} {#1}
	\addtocontents{toc}{\protect\setcounter{tocdepth}{2}} }


\newcommand{\bfone}{\mathbf1}
\newcommand{\ift}{\infty}
\newcommand{\q}{\mbox{ }}
\newcommand{\qd}{\quad}


\vfuzz2pt % Don't report over-full v-boxes if over-edge is small
\hfuzz2pt % Don't report over-full h-boxes if over-edge is small


\newcommand{\Mag}[1]{\begingroup\color{magenta} #1\endgroup} %Magenta 
\newcommand{\Bl}[1]{\begingroup\color{blue} #1\endgroup} %Blue

\newcommand{\ken}[1]{\textcolor{magenta}{[Ken: #1]}}
\newcommand{\kyu}[1]{\textcolor{blue}{[Kyudong: #1]}}
\newcommand{\you}[1]{\textcolor{cyan}{[Young-Jin: #1]}}
\newcommand{\inj}[1]{\textcolor{red}{[In-Jee: #1]}}
\newcommand{\kwan}[1]{\textcolor{orange}{[Kwan: #1]}}





\begin{document}
\title{Existence and stability of Sadovskii vortices:\\
from patch to smooth vortices}
\author{Ken Abe}
\address{Department of Mathematics, Graduate School of Science, Osaka Metropolitan University, 3-3-138 Sugimoto, Sumiyoshi-ku Osaka, 558-8585, Japan.}
\email{kabe@omu.ac.jp}
\author{Kyudong Choi}
\address{Department of Mathematical Sciences, Ulsan National Institute of Science and Technology, 50 UNIST-gil, Eonyang-eup, Ulju-gun, Ulsan 44919, Republic of Korea.}
\email{kchoi@unist.ac.kr}
\author{In-Jee Jeong}
\address{School of Mathematics, Korea Institute for Advanced Study, 85 Heogi-ro, Seoul 02455, Republic of Korea.}
\email{ijeong@kias.re.kr}
\author{Young-Jin Sim}
\address{School of Mathematics, Korea Institute for Advanced Study, 85 Heogi-ro, Seoul 02455, Republic of Korea.}
\email{yjsim@kias.re.kr}
\author{Kwan Woo}
\address{Departement Mathematik und Informatik, Universit\"at Basel, CH-4051 Basel, Switzerland.}
\email{kwan.woo@unibas.ch}
\date{Source: arXiv:2507.00910v3; distributed under CC BY 4.0}
\renewcommand{\thefootnote}{\fnsymbol{footnote}}
\footnotetext{\emph{Key words: two-dimensional Euler equations, Sadovskii vortex, orbital stability, variational method, sharp energy inequality, concentration--compactness} \\
\emph{2020 AMS Mathematics Subject Classification: 35Q31, 76B47, 35A15}}
\renewcommand{\thefootnote}{\arabic{footnote}}
	\begin{abstract}  
    We establish a scaling-invariant variational framework for steadily translating dipoles of the two-dimensional incompressible Euler equations. Specifically, we consider the maximization of the kinetic energy subject to constraints on the impulse and the $L^p$-norm ($1<p\leq \infty$) of vorticity without imposing any restriction on the total mass. In contrast to variational constructions with a mass constraint, the shape of the vortices in our framework is determined directly by the scaling-invariant structure of the energy.

    We prove that every vortex arising from this variational principle necessarily touches the symmetry axis and is therefore what is known in the physics literature as a \emph{Sadovskii vortex}, a configuration known to emerge as the endpoint of steady vortex-dipole branches. The construction is based on a sharp energy inequality under the two constraints of fixed impulse and $L^p$-norm, combined with an application of the concentration--compactness principle. This yields a family of axis-touching solutions parameterized by $p \in (1,\infty]$, interpolating between the classical Sadovskii vortex patch (corresponding to $p=\infty$) and vortices of arbitrarily high regularity as $p \to 1$. In particular, the classical Chaplygin--Lamb dipole (corresponding to $p=2$) is recovered within this family.

    Previous variational constructions treated only particular cases, notably $p=2$ and $p=\infty$, and were formulated under an additional mass constraint. Under this mass constraint, the natural scaling leads to the restriction $p>4/3$. By removing the mass constraint, we obtain a unified scaling-invariant variational principle valid for all $1<p\le\infty$. The endpoint $p=1$ nonetheless remains a natural limiting threshold for the present framework.

    As a consequence of the variational structure, we establish {an orbital} stability result, demonstrating that the axis-touching geometry persists under small perturbations. Finally, we derive a quantitative bound on the horizontal center of mass of perturbed solutions, showing that they propagate at nearly the same speed as the underlying Sadovskii vortex.
	\end{abstract}

\maketitle
\noindent\emph{Source, version and licence.} \emph{Existence and stability of Sadovskii vortices: from patch to smooth vortices}, by Ken Abe, Kyudong Choi, In-Jee Jeong, Young-Jin Sim and Kwan Woo. The source of this editable unit is the e-print of \textbf{version v3} of arXiv:2507.00910, https://arxiv.org/abs/2507.00910v3, whose material is distributed under the Creative Commons Attribution 4.0 International licence, http://creativecommons.org/licenses/by/4.0/, as recorded in the arXiv licence metadata for this paper and shown on that abstract page. The whole of the body below is version v3, the newest version of the paper. Full rights notice: licenses/arxiv-2507.00910-CC-BY-4.0.txt.\par\smallskip
\noindent\textit{Editorial scope.} Counted unit: original Theorem 6.1 of the article, the statement carried by the source environment \texttt{theorem} with source label \texttt{thm\_main\_1}, cached-originals/source0.tex lines 2146--2190, which the article states in the same words as Theorem 1.2 in the introduction (source lines 509--529, delivered here as uncounted context) and which carries the optional title ``Theorem 1.2: Sadovskii vortex family''. It is delivered together with the author's own complete proof as the single primary proof body, source0.tex lines 2192--2271 (80 lines): the \texttt{proof} environment that immediately follows the statement, that disposes of existence, of the Sadovskii-vortex geometry and of the regularity of part (ii) through the named results of the same article, and that then proves the orbital stability of part (iii) in full, first for $1<p<\infty$ by the compactness-contradiction argument and then for $p=\infty$ by the decomposition and energy estimates, before it closes. No proof marker is added and no mathematics is written, repaired or re-derived; the author's notation and the printed slips of the source are kept literal. Necessary complete context is retained uncounted, in source order: source0.tex lines 402--410, the author's own abstract; lines 445--453, the half-plane Euler equation to which the stability statement refers; lines 461--484, Definition 1.1 of a Sadovskii vortex; lines 509--529, the introduction statement of the same theorem; lines 647--654, the half Euler equation of the function space setting; lines 1098--1116, the energy identity; lines 1205--1208 and 1256--1258, the maximization problem and the two-parameter scaling; lines 1333--1341 and 1666--1684, the auxiliary estimates and the Euler--Lagrange identities the proof invokes; lines 1699--1705, 1722--1725, 1732--1741, 1786--1789, 1860--1863 and 1896--1909, Lemma 4.1 on regularity with Remark 4.2, Lemma 4.3 (Pohozaev identities), Lemma 4.5 (compact support), the touching lemma and the optimal-regularity remark; and lines 1940--1959, Theorem 5.2 on the existence of maximizers, which is the compactness input of the counted proof. The article's own numbering is reproduced by explicit counter settings: the counter \texttt{theorem}, declared by \texttt{\textbackslash newtheorem\{theorem\}\{Theorem\}[section]} and shared by corollary, lemma, proposition, definition, example and remark, and the section-scoped equation counter, are set before each retained passage, so the counted statement prints as Theorem 6.1 with its subtitle ``(Theorem 1.2)'', Definition 1.1 as 1.1, Theorem 1.2 as 1.2, Lemma 4.1 as 4.1, Remark 4.2 as 4.2, Lemma 4.3 as 4.3, Lemma 4.5 as 4.5, Theorem 5.2 as 5.2, and the retained displays keep the tags they carry in the article. Every \texttt{\textbackslash ref} and \texttt{\textbackslash eqref} inside every retained passage resolves inside this delivered file, and no \texttt{\textbackslash cite} occurs in any retained passage: the closure of the retained passages over references and citations was computed and contains no citation at all, so this unit delivers no bibliography entry and the article's own bibliography data file \texttt{Sadovskii.bib} is neither spliced into the bundled source nor redistributed. The e-print ships one artwork file, \texttt{center\_speed.pdf}, used by the single figure environment at source0.tex lines 797--804 with an \texttt{\textbackslash includegraphics} command; that figure lies outside every retained passage, so no figure environment and no graphics-inclusion command occurs anywhere in this delivered file and no artwork is redistributed. The source's own class is AMS \texttt{amsart} (\texttt{\textbackslash documentclass[11pt]\{amsart\}}); the e-print ships no class or style file, so no publisher class is shipped, used or redistributed, and the wrapper is the same AMS \texttt{amsart} class with the paper's own macros and theorem declarations transcribed from its preamble. Every declared body of this unit is an exact, unmodified span of source0.tex: no body is a bounded passage comparison, \texttt{omit} was not used anywhere, and consequently no fidelity receipt is asserted in delivery-evidence.json. The complete concatenated original source is bundled as sources/180847/source0.tex. Omitted from this editable unit and therefore absent from it are source0.tex lines 1--444; 454--460; 485--508; 530--646; 655--1097; 1117--1204; 1209--1255; 1259--1332; 1342--1665; 1675--1675; 1685--1698; 1706--1721; 1726--1731; 1742--1785; 1790--1859; 1864--1895; 1910--1939; 1960--2145; 2191--2191; 2272--2832; among them the article's own preamble, title block and table of contents, Sections 2 and 3 in all but the retained estimates, most of Sections 4 and 5, the intro-level stability theorem (Theorem 1.3), the whole minimal-mass section with its shift estimates, the closing remarks, and the article's bibliography data. The AMS \texttt{amsart} wrapper, the named format packages of the pinned TeX Live runtime and this editorial source, licence and scope paragraph are editorial formatting only.\par\smallskip
\setcounter{section}{1}
\setcounter{equation}{0}
\setcounter{theorem}{0}
	\begin{equation}\label{eq: Euler eq.}
		\left\{
		\begin{aligned}
			&\partial_{t} \omega + \bfv \cdot \nabla \omega = 0,\\
			&\bfv=\nabla^\perp(-\Delta_{\bbR^2})^{-1}\omega,\\
			&\omega|_{t=0} = \omega_{0},
		\end{aligned}
		\right.
	\end{equation} where $\nabla^\perp = (\partial_{x_2}, -\partial_{x_1})$, and $\omg_0:\bbR^2\to\bbR$ is the  given initial data. The \textit{stream function} $\psi=\psi[\omg]:=(-\Delta_{\bbR^2})^{-1}\omega$ allows us to recover the \textit{velocity} $\bfv$ from the \textit{vorticity} $\omega$ through the Biot--Savart law:

\setcounter{section}{1}
\setcounter{equation}{1}
\setcounter{theorem}{0}
	\begin{definition}\label{def: Sadovskii vortex}
		A function $\overline{\omega} \in L^\infty(\bbR^2)$ is called a \textit{Sadovskii vortex} if it is \textbf{compactly supported} and satisfies the following properties:
		\begin{itemize}[itemsep=1em, topsep=1em]
			\item[(i)] \textbf{Traveling wave.} 
			There exists a constant $W > 0$ such that  
			\[
			\omega(t, \bfx) := \overline{\omega}(\bfx - (W, 0)t)
			\]
			solves the Euler equations \eqref{eq: Euler eq.}.  
			We refer to $W$ as the traveling speed of $\overline{\omega}$.
			
			\item[(ii)] \textbf{Odd symmetric dipole.}
			The vorticity $\overline{\omega}$ is odd in the $x_2$-variable:
			\[
			\overline{\omega}(x_1, x_2) = -\overline{\omega}(x_1, -x_2) \geq 0 
			\quad \mbox{for }\, x_1 \in \bbR, \,\, x_2 > 0.
			\]
			
			\item[(iii)] \textbf{Touching.} The dipole exhibits full contact; there exist $r > 0$ and $c \in \bbR$ such that  
			\begin{equation}\label{touching_scale}
			\left\{ \bfx \in \bbR^2 : |\bfx - (c, 0)| < r \right\} \subset \supp \overline{\omega}.
			\end{equation}
		\end{itemize}
	\end{definition}

\setcounter{section}{1}
\setcounter{equation}{2}
\setcounter{theorem}{1}
	\begin{theorem}[Theorem~\ref{thm_main_1}: Sadovskii vortex family]\label{thm_intro_main}
		For each $p \in (1,\infty]$, there exists a nonempty set $\mathcal{S}(p)$ of vorticities defined on $\bbR^2_+$ satisfying the following:
		\begin{itemize}[itemsep=0.8em, topsep=0.8em]
			
			\item[(i)] \textbf{Sadovskii vortex.}  	 The odd extension of every element in $\mathcal{S}(p)$ is a Sadovskii vortex with traveling speed $W=W(p)$ and touching scale $r=r(p)$ satisfying \eqref{touching_scale} in Definition~\ref{def: Sadovskii vortex}.

			
			
			\item[(ii)] \textbf{Regularity transition: from patch to smooth vortices.}  
			Every element of $\mathcal{S}(\infty)$ is a vortex patch, i.e. for each $\omg \in \mathcal{S}(\infty)$, there is $A \subset \mathbb{R}^{2}_{+}$ such that $\omg = \mathbf{1}_{A}$.
			As $p \searrow 1$, the regularity of the corresponding Sadovskii vortices increases. In particular, for every $m \in \mathbb{N}$, there exist Sadovskii vortices of class $C^m(\mathbb{R}^2)$.

			
			\item[(iii)] \textbf{{Orbital} stability.}  
			For any $1 < p < \infty$, the set $\mathcal{S}(p)$ is {orbitally} stable with respect to the norm  $\|\cdot\|_{L^p} + \|x_2(\cdot)\|_{L^1}$.
            {
            For $p = \infty$, stability in the norm $\|x_2(\cdot)\|_{L^1}$ holds under the smallness assumption for the excess $L^q$-norm $1 < q < \infty$ on the initial data, which is specified in Theorem~\ref{thm_main_1}-(iii).
            }
		\end{itemize}
		More precisely, each $\omega \in \mathcal{S}(p)$ maximizes the kinetic energy under unit impulse and unit $L^p$ norm constraints, subject to positivity on $\bbR^2_+$ and odd-symmetry, with no restriction on the total mass.
	\end{theorem}

\setcounter{section}{1}
\setcounter{equation}{4}
\setcounter{theorem}{5}
	\begin{equation}\label{eq: half Euler eq.}
		\left\{
		\begin{aligned}
			& \partial_{t}\omega + \bfv \cdot \nabla\omg = 0 \quad \mbox{ in } \,  [0,\infty)\times\bbR^2_+,\\
			& \bfv =  \nabla^\perp(-\Delta_{\bbR^2_+})^{-1}\omega \quad \mbox{ in }\, [0,\infty)\times\bbR^2_+,
		\end{aligned}
		\right.
	\end{equation} where the stream function $\calG[\omg]:=(-\Delta_{\bbR^2_+})^{-1}\omg$ is given by

\setcounter{section}{2}
\setcounter{equation}{9}
\setcounter{theorem}{4}
		\begin{align}
			&\iint_{\bbR^2_+ \times \bbR^2_+} G(\bfx,\bfy)\,|\omega(\bfx)\,\eta(\bfy)|\,\dd\bfx\,\dd\bfy
			\, 
			\lesssim_p \, 
			\|x_2\omega\|_{1}^{\frac{2p-2}{3p-2}}\,
			\|\omega\|_{p}^{\frac{p}{3p-2}}\,
			\|x_2\eta\|_{1}^{\frac{2p-2}{3p-2}}\,
			\|\eta\|_{p}^{\frac{p}{3p-2}},   \label{eq_0704_01}\\
			&E(\omega) 
			\lesssim_p 
			\|x_2\omega\|_{1}^{\frac{4p-4}{3p-2}}\,
			\|\omega\|_{p}^{\frac{2p}{3p-2}},		\label{eq_0704_02}  \\
			&|E(\omega)-E(\eta)| \;
			\lesssim_p\;
			\|x_2(\omega+\eta)\|_{1}^{\frac{2p-2}{3p-2}}\,
			\|\omega+\eta\|_{p}^{\frac{p}{3p-2}}\,
			\|x_2(\omega-\eta)\|_{1}^{\frac{2p-2}{3p-2}}\,
			\|\omega-\eta\|_{p}^{\frac{p}{3p-2}}.  \label{eq_0704_03}
		\end{align}

\setcounter{section}{2}
\setcounter{equation}{12}
\setcounter{theorem}{6}
		\begin{align}
			\sup_{|\bfx| \ge R}\calG[\omega](\bfx)&\to 0 \quad \text{as } \,\, R\to\infty, \label{eq_vanishing} \\
			E(\omega)&=\frac{1}{2}\left\|\nabla \calG[\omega] \right\|_{2}^{2}.  \label{eq_L2energy}
		\end{align}

\setcounter{section}{2}
\setcounter{equation}{16}
\setcounter{theorem}{6}
	\begin{equation}\label{two_scaling}
		\overline{\omega}(\bfx) = \frac{\tau^3}{\mu} \omega(\tau \bfx),\quad \tau=\left(\frac{\mu}{\sigma}\right)^{\frac{p}{3p-2}},
	\end{equation}

\setcounter{section}{3}
\setcounter{equation}{0}
\setcounter{theorem}{0}
	\begin{equation}
		\left\{
		\begin{aligned}
			\kappa \, \omega(\bfx)^{p-1}&=(\mathcal{G}[\omega](\bfx)-Wx_2)_{+},\quad 1<p<\infty,\\
			\omega(\bfx)&=\sigma \mathbf{1}_{\{\bfy\in\mathbb{R}^2_+\,:\,\mathcal{G}[\omega](\bfy)-Wy_2>0 \}}(\mathbf{x}),\quad p=\infty, 
		\end{aligned}
		\right.
		\label{eq: EL}
	\end{equation}

\setcounter{section}{4}
\setcounter{equation}{0}
\setcounter{theorem}{0}
			\begin{equation}
				\left\{
				\begin{aligned}
					-\Delta \psi (\bfx)&= f\left(\psi \left(\bfx \right) - Wx_2  \right)\quad \textrm{ in }\, \bbR^2_+, \\
					\psi&=0 \quad \textrm{ on }\, \partial\bbR^2_+,
				\end{aligned}
				\right.
				\label{eq_ell}
			\end{equation}

\setcounter{section}{4}
\setcounter{equation}{1}
\setcounter{theorem}{0}
			\begin{equation}    
				\begin{aligned}
					f(s) = \begin{cases}
						\kappa^{-\frac{1}{p-1}} s_+^{\frac{1}{p-1}}, & \text{for }\,\, 1<p<\infty, \\[5pt]
						\mathbf{1}_{\{s > 0\}}, & \text{for}\,\, p = \infty.
					\end{cases}
				\end{aligned}
				\label{f}
			\end{equation}

\setcounter{section}{4}
\setcounter{equation}{2}
\setcounter{theorem}{0}
			\begin{lemma}\label{lem_regularity_1}
				Let $1<p \leq \infty$ and $\omega \in \calS_{1,1}(p)$. Then, $\omega\in L^{\infty}(\bbR^2_+)$ and $\psi =\calG[\omega] \in L^{\infty}(\mathbb{R}^{2}_{+})\cap W^{2,q}_{\mathrm{loc}}(\overline{\mathbb{R}_{+}^{2}})$ for all $q\in (1,\infty)$. In particular, $\psi\in C^{1,\beta}(\overline{\mathbb{R}^{2}_{+}})$ for all $\beta \in (0, 1)$. For $p\in(1,\infty)$, $\omega\in C^{m, \alpha}(\overline{\bbR^2_{+}})$ and 
				\begin{align}
					m = \left\lceil \frac{1}{p-1} - 1 \right\rceil, \quad
					\alpha = \frac{1}{p-1} - m \in (0,1]. \label{malpha}
				\end{align}
			\end{lemma}

\setcounter{section}{4}
\setcounter{equation}{3}
\setcounter{theorem}{1}
			\begin{remark}\label{rmk_reg_odd}
				The regularity stated in Lemma~\ref{lem_regularity_1} is preserved under odd extension across the boundary $\{x_2 = 0\}$.
				This is a consequence of the Euler--Lagrange equation \eqref{eq_ell} and the degeneracy of the nonlinearity $f$ in \eqref{f}, whose derivatives up to order $m = \left\lceil \frac{1}{p-1} - 1 \right\rceil$ vanish at the origin.
			\end{remark}

\setcounter{section}{4}
\setcounter{equation}{3}
\setcounter{theorem}{2}
			\begin{lemma}
				\label{lem_pohozaev}
				Let $1<p< \infty$ and $\omega \in \calS_{1, 1}(p)$.
				The constants $\kappa_p  > 0$ and $W_p> 0$ in \eqref{eq: EL} satisfy 
				\begin{equation}
					\label{eq_pohozaev}
					W_p=\frac{4p-4}{3p-2}\calI_{1,1}(p),\quad \kappa_p =\frac{2p}{3p-2} \calI_{1,1}(p).
				\end{equation}
				For $p=\infty$, $W_{\infty}=\frac{4}{3}\calI_{1,1}(\infty)$.
			\end{lemma}

\setcounter{section}{4}
\setcounter{equation}{4}
\setcounter{theorem}{4}
			\begin{lemma}
				\label{lem_cpt_support}
				Let $1<p\leq \infty$ and $\omega \in \calS_{1,1}(p)$. Then, $\supp \omega = \overline{ \{ \bfx \in \bbR^2_+: \omega > 0 \} }$ is compact.
			\end{lemma}

\setcounter{section}{4}
\setcounter{equation}{7}
\setcounter{theorem}{8}
			\begin{lemma}
				\label{prop_touching}    
				Let $1<p \leq \infty$. There exists a constant $r = r(p) > 0$ such that $B_{r}(\mathbf{0}) \subset \supp \widetilde{\omega}$ for $\omega \in \calS_{1,1}({p})$ by a suitable translation in the $x_1$-direction.
			\end{lemma}

\setcounter{section}{4}
\setcounter{equation}{7}
\setcounter{theorem}{9}
			\begin{remark}
				\label{rmk_optimal_regularity}
				The regularity result ${\omega} \in C^{m, \alpha}(\overline{\bbR^2_{+}})$ for $1<p< \infty$ in Lemma \ref{lem_regularity_1} is optimal in the sense that ${\omega} \notin C^{m, \alpha+\varepsilon}(\overline{\bbR^2_{+}})$ for any $\varepsilon > 0$.
				Indeed, for $\omega \in \calS_{1,1}(p)$, $\psi=\mathcal{G}[\omega] \in C^2(\overline{\mathbb{R}^2_+})$ satisfies $\psi(x_1, 0) = 0$ for $x_1 \in \mathbb{R}$ and  
				\[
				\psi(0, x_2)-Wx_2 = (\partial_{x_2}\psi(0,0)-W)x_2 + o(x_2) \quad \textrm{as}\,\, x_2\to0.
				\]
				This means that 
				\[
				\kappa^{p-1} \omega(0, x_2) = \left( (\partial_{x_2}\psi(0,0) - W)x_2+ o(x_2)\right)^{\frac{1}{p - 1}}_+ 
				\ \simeq \ (\partial_{x_2}\psi(0,0) - W)^{\frac{1}{p - 1}} \, x_2^{\frac{1}{p - 1}},
				\]
				and ${\omega}\notin C^{m, \alpha+ \varepsilon}(\overline{\bbR^2_{+}})$ for any $\varepsilon > 0$. 
			\end{remark}

\setcounter{section}{5}
\setcounter{equation}{0}
\setcounter{theorem}{1}
			\begin{theorem}\label{thm_existence_new}
				Let $1<p\leq \infty$ and $0<\mu, \sigma <\infty$. For any non-negative sequence $\{\omega_n\}_{n}$ on $\bbR^{2}_{+}$  satisfying 
				\begin{equation}\label{eq_stab_energy}
					\begin{aligned}
						{\liminf_{n \to \infty} E(\omega_n) \ge \calI_{\mu, \sigma}(p), \quad
							\limsup_{n \to \infty}\|x_2 \omega_n\|_1} \le \mu, \quad
						\limsup_{n \to \infty} \|\omega_n\|_p \le \sigma, 
					\end{aligned}
				\end{equation}
				there exists $\omega \in \calS_{\mu, \sigma}(p)$ such that, by a suitable translation for the $x_1$-variable, $\{\omega_{n} \}_n$ subsequently converges to $\omega$ with the norm $||x_2 (\cdot)||_{1}$ and 
				
				\begin{equation}
					\left\{
					\begin{aligned}
						\omega_{n}&\to \omega\quad \textrm{in}\ L^{p}(\mathbb{R}^{2}_{+})\quad \textrm{ for }\ p \in (1, \infty),\\
						\omega_{n} &\rightharpoonup^* \omega \quad \textrm{in}\ L^{\infty}(\mathbb{R}^{2}_{+})\quad \textrm{ for }\ p=\infty.
					\end{aligned}
					\right.
				\end{equation}
			\end{theorem}

\setcounter{section}{6}
\setcounter{equation}{0}
\setcounter{theorem}{0}
			\begin{theorem}[Theorem~\ref{thm_intro_main}]
				\label{thm_main_1}
				Let $1<p \leq \infty$ and $0<\mu,\sigma<\infty$. The following holds:
				
				\begin{enumerate}[itemsep=1em, topsep=1em]
					\item[(i)] \textbf{(Existence)} The odd extension $\tld{\omega}_p$ of each $\omega_p\in \calS_{\mu,\sigma}(p)$ is a Sadovskii vortex travelling with the same speed $W = W(p, \mu, \sigma) > 0$ and 
					\begin{equation}\label{eq_main_supp}
						B_{r}(\mathbf{0}) \subset \supp \tld{\omega}_p,
					\end{equation}
					with the uniform constant $ r = r(p, \mu, \sigma) >0$ by a suitable translation for the $x_1$-variable. 
					\item[(ii)] \textbf{(Regularity)} For $p=\infty$, $\tld{\omega}_{\infty}$ is a Sadovskii vortex patch. For $p \in (1, \infty)$, 
					\begin{equation}\label{eq_main_regularity}
						\tld{\omega}_p \in C^{m, \alpha}(\bbR^2), \quad \textrm{ for } \, m =m_p = \left\lceil \frac{1}{p-1} - 1 \right\rceil,\ \alpha = \alpha_p =\frac{1}{p-1}-m,
					\end{equation}
					and $\tld{\omega}_p \notin C^{m, \alpha + \varepsilon}(\bbR^2)$ for any $\varepsilon > 0$.
					\item[(iii)] \textbf{(Stability)} For each $p \in (1, \infty)$ and $\varepsilon>0$, there exists $\delta = \delta(p, {\mu, \sigma,\, }   \varepsilon)>0$ such that for any initial data $0 \le \theta_0 \in  L_c^{\infty}(\bbR^2_+)$ satisfying  
					\begin{align}
						\inf_{\omega\in \calS_{\mu,\sigma}(p)}\left\{  \left\| \theta_0-\omega \right\|_{p} 
						+\left\| x_2 \left( \theta_0-\omega \right) \right\|_{1} \right\}   \label{eq: initial}
						\leq  \delta,
					\end{align}
					the unique global-in-time solution $\theta(t)$ to 
                    {\eqref{eq: half Euler eq.}}
                    %\eqref{eq: Euler eq.} 
                    satisfies
					\begin{align} 
						\inf_{\omega\in \calS_{\mu,\sigma}(p)}\left\{   \left\| \theta \left( t \right)-\omega \right\|_{p}
						+\left\|x_2 \left( \theta \left(t \right)-\omega  \right) \right\|_{1}\right\}
						\leq  \varepsilon \quad \mbox{for all } \,\, t \in \bbR.   \label{eq: stability}
					\end{align} 
                   
                    For $p = \infty$, $\varep >0$, and each fixed $q \in (1, \infty)$, there exists $\delta = \delta(\mu, \sigma, \varepsilon, q) > 0$ such that for any initial data $0 \le \theta_0 \in  L_c^{\infty}(\bbR^2_+)$ satisfying
                    \[
                     {\| (\theta_0 - \sigma)_+ \|_q}  +   \inf_{\omega \in \calS_{\mu, \sigma}(\infty)} \| x_2 (\theta_0 - \omega ) \|_1 \le \delta,
                    \]
                    the corresponding solution $\theta(t)$ to 
                    {\eqref{eq: half Euler eq.}}
                    %\eqref{eq: Euler eq.} 
                    satisfies
                    \[
                    \inf_{\omega \in \calS_{\mu, \sigma}(\infty)} \| x_2(\theta(t) - \omega) \|_1 \le \varepsilon \quad \textrm{for all } \,\, t \in \bbR.
                    \]
                    
				\end{enumerate}
			\end{theorem}

\setcounter{section}{6}
\setcounter{equation}{4}
\setcounter{theorem}{1}
			\begin{proof}
				The existence of maximizers $\omega_p\in \calS_{\mu,\sigma}(p)$ follows from Theorem \ref{thm_existence_new}.
                By the Euler--Lagrange equation \eqref{eq: EL} and Lemma~\ref{lem_pohozaev}, together with two-parameter scaling \eqref{two_scaling}, any maximizer $\omega_p \in \calS_{\mu,\sigma}(p)$ is a traveling wave with traveling speed $W = W(p, \mu, \sigma)$.
				The support of $\omega_p$ is compact by Lemma~\ref{lem_cpt_support}, and in contiguous contact with the $x_1$-axis by Lemma \ref{prop_touching}, with touching scale $r = r(p, \mu, \sigma)$ via the scaling \eqref{two_scaling}.
                Thus, $\tld{\omega}_p$ is the Sadovskii vortex.
                The regularity of $\tld{\omega}_p$ follows from Lemma \ref{lem_regularity_1} and Remarks~\ref{rmk_reg_odd} and \ref{rmk_optimal_regularity}.
				
				It remains to show the stability. For $1<p<\infty$, suppose that the claim were false. Then there exist $\varepsilon_0>0$, non-negative $\{\theta_{0, n} \}_{n = 1}^{\infty}\subset L^{\infty}_{c}(\mathbb{R}^{2}_{+})$, and $\{t_n\}_{n = 1}^{\infty}\subset \mathbb{R}$ such that for each $n \ge 1$, 
				\begin{align*}\label{eq241228_2}
					&\inf_{\omega \in \calS_{\mu,\sigma}(p)} \left\{  \left\| \theta_{0, n} - \omega \right\|_{p} + \left\| x_2 (\theta_{0, n} - \omega) \right\|_{1} \right\}   \le  \frac{1}{n},\\
					&\inf_{\omega \in \calS_{\mu,\sigma}(p)}\left\{  \left\| \theta_{n} \left(t_n, \cdot \right) - \omg \left( \cdot  \right) \right\|_{p} + \left\| x_2 \big( \theta_{n} \left(t_n, \cdot \right) - \omega \left( \cdot  \right) \big) \, \right\|_{1}  \right\}  \ge  \varepsilon_0>0,
				\end{align*}
				for the global-in-time weak solution $ \theta_n(t)$ to \eqref{eq: Euler eq.} satisfying $\theta_n(0) = \theta_{0, n}$. We take a sequence $\{\omega_n\}_{n} \subset \calS_{\mu,\sigma}(p)$ such that
				\[
				\| x_2 (\theta_{0, n} - \omega_n) \|_{1} + \| \theta_{0, n} - \omega_n\|_{p}  \to  0,
				\]
				and observe from \eqref{eq_0704_03} that  
				\[
				\lim_{n \to \infty}E(\theta_{0, n}) = \calI_{\mu,\sigma}(p),\quad \limsup_{n\to\infty}||x_2\theta_{0, n}||_{1}\leq \mu, \quad \limsup_{n\to\infty}||\theta_{0, n} ||_{p}\leq \sigma.
				\]
				By the conservation property of the Yudovich solution, $\set{\theta_{n}(t_n)}_{n=1}^\infty$ also satisfies the same property. By applying Theorem~\ref{thm_existence_new} to the sequence $\{ \theta_n(t_n)  \}_{n}$, we obtain a subsequence still denoted by $\{ \theta_n(t_n)  \}_{n}$ such that $\theta_n(t_n)$ converges to some ${\omega} \in \calS_{\mu,\sigma}(p)$ with the norm $||\cdot ||_{p}+||x_2(\cdot)||_{1}$ by a suitable translation for the $x_1$-variable. By letting $n \to \infty$,
				\begin{align*}
					0 &=  \lim_{n \to \infty} \left[	\left\| x_2 \left( \theta_n (t_n) - {\omega} \right) \right\|_{1} +    \left\|  \theta_n (t_n) - {\omega} \right\|_{p} \right] \\
					&\ge  \liminf_{n \to \infty} \left\{ \inf_{\omega \in \calS_{\mu,\sigma}(p)} \left(   \left\| x_2 \left( \theta_{n} \left(t_n\right) - \omega \right)\, \right\|_{1} + \left\| \theta_{n} \left(t_n \right) - \omega \right\|_{p} \right) \right\}  \ge  \varepsilon_0>0.
				\end{align*}
				We obtain a contradiction.
				
                {
                For $p = \infty$, fix $q \in (1, \infty)$ and suppose that the claim were false.
                Then there exist $\varepsilon_0 > 0$, non-negative $\theta_{0,n} \in L^\infty_c(\mathbb{R}^2_+)$, times $t_n \in \bbR$, and $\omega_n \in \calS_{\mu,\sigma}(\infty)$ such that
                \[
                \|x_2(\theta_{0,n} - \omega_n)\|_1 + \|(\theta_{0,n} - \sigma)_+\|_q \to 0,
                \]
                whereas the corresponding solutions satisfy
                \[
                \inf_{\omega \in \calS_{\mu,\sigma}(\infty)} \|x_2(\theta_n(t_n) - \omega)\|_1 \ge \varepsilon_0.
                \]
                Set
                \[
                v_n(t) = \theta_n(t) \wedge \sigma, \quad \eta_n(t) = (\theta_n(t) - \sigma)_+.
                \]
                Conservation of impulse and of the vorticity distribution gives
                \[
                \|x_2\theta_n(t)\|_1 = \|x_2\theta_{0,n}\|_1 \to \mu, \quad \|\eta_n(t)\|_q = \|\eta_n(0)\|_q \to 0.
                \]
                Since $0 \le v_n(t) \le \sigma$ and $0 \le \eta_n(t) \le \theta_n(t)$, \eqref{eq_0704_02} yields, uniformly in $t$,
                \[
                E(v_n(t)) \le C, \quad E(\eta_n(t)) \le C_q \|x_2\theta_{0,n}\|_1^{\frac{4q-4}{3q-2}} \|\eta_n(0)\|_q^{\frac{2q}{3q-2}} \to 0.
                \]
                By \eqref{eq_L2energy} and the Cauchy--Schwarz inequality,
                \[
                0 \le E(\theta_n(t)) - E(v_n(t)) \le 2\sqrt{E(v_n(t)) E(\eta_n(t))} + E(\eta_n(t))\to 0
                \]
                uniformly in $t$.
                Since $0\le \omega_n \le \sigma$,
                \[
                \|x_2(v_n(0) - \omega_n)\|_1 \le \|x_2(\theta_{0,n} - \omega_n)\|_1 \to 0.
                \]
                Then, by applying \eqref{eq_0704_03} with $p=\infty$ we have $E(v_n(0)) \to \calI_{\mu,\sigma}(\infty)$.
                Hence energy conservation and the preceding estimates imply
                \[
                E(v_n(t_n))\to\calI_{\mu,\sigma}(\infty), \quad \limsup_{n\to\infty}\|x_2v_n(t_n)\|_1 \le \mu, \quad \|v_n(t_n)\|_\infty \le \sigma.
                \]
                By Theorem~\ref{thm_existence_new}, after passing to a subsequence and applying horizontal translations,
                \[
                \|x_2(v_n(t_n) - \omega) \|_1 \to 0
                \]
                for some $\omega \in \calS_{\mu,\sigma}(\infty)$.
                We apply the same translations to $\theta_n(t_n)$ and $\eta_n(t_n)$ and use the same notation.
                In particular, $\|x_2 v_n(t_n)\|_1 \to \mu$, and therefore
                \[
                \|x_2\eta_n(t_n)\|_1 = \|x_2\theta_{0,n}\|_1 - \|x_2 v_n(t_n)\|_1 \to 0.
                \]
                Consequently,
                \[
                \|x_2(\theta_n(t_n)-\omega)\|_1 \le \|x_2(v_n(t_n)-\omega)\|_1 + \|x_2\eta_n(t_n)\|_1 \to 0.
                \]
                Since $\calS_{\mu,\sigma}(\infty)$ is invariant under horizontal translations, this contradicts the choice of $t_n$.
                }
			\end{proof}

\end{document}
